% !TeX root = ../main.tex
\section{Introduction}
\label{sec:intro}

Guitarists learn songs from tablature: a notation that gives, for every note,
the string to play and the fret to press, together with rhythm and marks for
techniques such as bends, slides, hammer-ons and vibrato. Tablatures for
popular music are mostly written by hand, by ear, and shared on community
sites with uneven quality. Writing them automatically from the recording
itself is therefore a long-standing goal of music information retrieval (MIR),
known as guitar tablature transcription (GTT).

GTT inherits every difficulty of automatic music transcription (AMT) and adds
its own. The guitar is polyphonic, and strummed chords place up to six notes
within a few tens of milliseconds. The same pitch can be played on up to five
string-fret positions, so pitch alone does not determine the tablature.
Electric guitar is almost always processed by amplifiers and effects whose
distortion creates intermodulation and dense harmonic spectra that confuse
pitch detectors. In commercial recordings the guitar shares the mix with
drums, bass, keyboards and voice, and is often double-tracked or layered with
other guitars. Finally, annotated data is scarce: labelling string, fret and
technique for every note requires either special hardware (hexaphonic pickups)
or expert transcription.

Two families of methods address the problem. \emph{Modular pipelines}
decompose it into sub-tasks (source separation, note transcription, beat and
bar tracking, string-fret assignment, technique detection) and solve each with
a dedicated model or algorithm. \emph{End-to-end} models predict string and
fret activity directly from the audio. Fig.~\ref{fig:families} contrasts the
two. A third line of work, large multi-instrument sequence-to-sequence
transcribers, does not produce tablature but has recently become the strongest
option for the note-transcription stage on full mixes, and we treat it as a
component of the modular family.

This review covers the task definition and its metrics
(Section~\ref{sec:task}), the datasets (Section~\ref{sec:datasets}), the
components of modular pipelines (Section~\ref{sec:modular}), end-to-end
audio-to-tablature models (Section~\ref{sec:e2e}), a benchmark of the models
with released weights under one protocol (Section~\ref{sec:benchmark}), and a
comparison with known failure modes and open problems
(Section~\ref{sec:comparison}). The scope is
the use case of transcribing a guitar part from a commercial song; results on
solo, clean recordings are reported, but read with that use case in mind.

\begin{figure}[t]
\centering
\begin{tikzpicture}[
  node distance=2.2mm,
  box/.style={draw, rounded corners=1pt, font=\scriptsize, align=center,
              minimum height=5mm, inner sep=2pt, fill=white},
  arr/.style={-{Stealth[length=1.6mm]}, semithick}]
  % Modular
  \node[box] (mix) {Mix};
  \node[box, right=of mix] (sep) {Source\\separation};
  \node[box, right=of sep] (amt) {Note\\transcription};
  \node[box, right=of amt] (fing) {String-fret\\assignment};
  \node[box, right=of fing] (tab) {Tab};
  \node[box, above=of amt] (beat) {Beats,\\bars};
  \node[box, below=of amt] (tech) {Techniques};
  \draw[arr] (mix) -- (sep);
  \draw[arr] (sep) -- (amt);
  \draw[arr] (amt) -- (fing);
  \draw[arr] (fing) -- (tab);
  \draw[arr] (mix.north) |- (beat.west);
  \draw[arr] (beat.east) -| (tab.north);
  \draw[arr] (sep.south) |- (tech.west);
  \draw[arr] (tech.east) -| (tab.south);
  \begin{scope}[on background layer]
    \node[draw, dashed, rounded corners, fit=(mix)(tab)(beat)(tech),
          inner sep=2.5mm, label={[font=\scriptsize]above left:Modular}] {};
  \end{scope}
  % End-to-end
  \node[box, below=13mm of mix] (mix2) {Solo\\guitar};
  \node[box, right=6mm of mix2, minimum width=36mm] (e2e)
       {Audio-to-tablature model\\(per-string frets, per frame or note)};
  \node[box, right=6mm of e2e] (tab2) {Tab};
  \draw[arr] (mix2) -- (e2e);
  \draw[arr] (e2e) -- (tab2);
  \begin{scope}[on background layer]
    \node[draw, dashed, rounded corners, fit=(mix2)(tab2), inner sep=2.5mm,
          label={[font=\scriptsize]below left:End-to-end}] {};
  \end{scope}
\end{tikzpicture}
\caption{The two families of tablature transcription systems. Modular
pipelines can start from a full mix; published end-to-end models start from
an isolated guitar recording.}
\label{fig:families}
\end{figure}
